\subsection{问题三：三类患者联合排程}

\subsubsection{背景需求的可评价定义}

门诊和体检附件包含开单与报告日期，却没有真实预约日期和历史预约时段。若直接读取历史报告时刻、机房和医生，排程将复现答案信息；若把报告日完全忽略，又无法定义门诊、体检服务水平。本文采用有限的回顾性压力场景：报告日期众数只确定计划日$\delta_i$，释放时刻不早于真实开单和计划日08:00，截止为计划日17:00：
\[
p_i^{B}=\max\{\mathrm{orderTime}_i,\delta_i+8\unit{h}\},\qquad
d_i^{B}=\delta_i+17\unit{h}.
\]
历史报告具体时刻、历史设备和医生均不进入调度。对背景事件集合$\I_B$，计划日代理完整完成率为
\[
R_B=\frac{\sum_{i\in\I_B}y_i}{|\I_B|}.
\]
该指标表示在标准班次下承接历史计划日需求的能力，不等于真实历史准点率。

\subsubsection{背景服务下限多目标模型}

问题三的两个主要目标为住院48小时率$R_I$与背景率$R_B$。超声预约与检查室指派适合在统一资源日历中评价\cite{chen2022}，等待保证也需要显式的服务等级约束\cite{bauerhenne2026}。加权和$\lambda R_I+(1-\lambda)R_B$对权重敏感，也会跳过非凸前沿，故本文采用$\varepsilon$约束：
\[
\begin{aligned}
\max\quad & \sum_{i\in\I_I}y_i,\\
\mathrm{s.t.}\quad
&R_B\ge\alpha,\\
&\text{问题二全部设备、医生、患者、准备与转运约束}.
\end{aligned}
\]
预先设置
\[
\alpha\in\{0,0.30,0.40,0.50,0.55,0.60,0.65,0.70\}.
\]
这里$\alpha$是全留出期背景完整完成率下限，与最终评价口径一致。算法同时记录逐日配额短缺天数，用于识别某些日期的项目结构不可行，但不以“每一天都恰好达到比例”替代全年服务约束。

\subsubsection{联合构造方法}

联合日历先装入3月30--31日的258个住院和286个背景预热事件。对留出期每个计划日，按训练期项目稀缺权重和释放时刻，从当日门诊、体检事件中选择至少
\[
Q_d(\alpha)=\left\lceil\alpha|\I_{B,d}|\right\rceil
\]
个背景目标事件，先尝试整体插入；随后使用问题二选定的共享FCFS处理所有已释放住院事件；最后再用剩余容量尝试当日未作为配额目标的背景事件。暖启动背景、住院和剩余背景均调用同一事务式患者级可行性内核，不存在三套互相矛盾的时长或能力规则。

这种次序不是把背景优先级永久置于住院之上。$\alpha$只预留最低背景服务；达到下限后，住院以48小时目标竞争容量；当天最后仍有空闲时，再继续安排背景。由于项目兼容和准备约束，有些背景事件即使在配额内也无法排入，故实际$R_B$通常高于$\alpha$，但不是机械相等。

\subsubsection{Pareto前沿与推荐点}

表\ref{tab:p3-pareto}列出八个运行点。$\alpha=0$时，住院先用容量，住院率为85.29\%，背景率仍可达到40.44\%；当$\alpha$提高到0.30，背景率升至59.53\%，住院率降至60.38\%。继续提高$\alpha$，背景率逐步升至65.19\%，住院率下降到57.01\%。所有满足$R_B\ge\alpha$的点均参与非支配判断；$\alpha=0.65$和0.70因实际背景率未达到目标，不作为该下限的可行解，但保留在表中以说明容量边界。

\begin{table}[H]
\centering
\caption{问题三背景下限网格及联合排程结果}
\label{tab:p3-pareto}
\resizebox{\textwidth}{!}{\begin{tabular}{rrrrrr}
\toprule
$\alpha$ & 背景完成数 & 背景率 & 住院完成数 & 住院率 & P90/h \\
\midrule
0.00 & 38,872 & 40.47\% & 60,368 & 85.30\% & 22.08 \\
0.05 & 38,872 & 40.47\% & 60,368 & 85.30\% & 22.08 \\
0.10 & 38,872 & 40.47\% & 60,368 & 85.30\% & 22.08 \\
0.15 & 38,872 & 40.47\% & 60,368 & 85.30\% & 22.08 \\
0.20 & 38,872 & 40.47\% & 60,368 & 85.30\% & 22.08 \\
0.25 & 38,872 & 40.47\% & 60,368 & 85.30\% & 22.08 \\
0.30 & 38,872 & 40.47\% & 60,368 & 85.30\% & 22.08 \\
0.35 & 38,872 & 40.47\% & 60,368 & 85.30\% & 22.08 \\
0.40 & 38,872 & 40.47\% & 60,368 & 85.30\% & 22.08 \\
\textbf{0.45} & \textbf{54,137} & \textbf{56.37\%} & \textbf{45,861} & \textbf{64.80\%} & \textbf{21.83} \\
0.50 & 55,708 & 58.00\% & 44,439 & 62.79\% & 22.00 \\
0.525 & 56,281 & 58.60\% & 43,852 & 61.96\% & 21.92 \\
0.55 & 57,186 & 59.54\% & 43,069 & 60.86\% & 22.33 \\
0.575 & 59,996 & 62.47\% & 42,440 & 59.97\% & 30.33 \\
0.60 & 61,075 & 63.59\% & 41,684 & 58.90\% & 30.42 \\
0.65 & 63,085 & 65.68\% & 40,273 & 56.91\% & 30.33 \\
\bottomrule
\end{tabular}
}
\end{table}

\begin{figure}[H]
\centering
\includegraphics[width=.78\textwidth]{fig05_pareto.pdf}
\caption{共享资源下的住院—背景服务Pareto前沿}
\label{fig:pareto}
\end{figure}

为避免主观挑选，定义可达前沿上的理想点$(R_B^{\max},R_I^{\max})$，对每个非支配点计算
\[
D(\alpha)=
\sqrt{
\left(\frac{R_B^{\max}-R_B(\alpha)}{R_B^{\max}-R_B^{\min}}\right)^2+
\left(\frac{R_I^{\max}-R_I(\alpha)}{R_I^{\max}-R_I^{\min}}\right)^2
}.
\]
选择距离最小者，并以背景率较高、住院率较高依次破同序。推荐点为$\alpha=0.30$：住院完整完成42\,730例，$R_I=60.38\%$；背景完整完成57\,177例，$R_B=59.53\%$；住院条件等待P50和P90分别为3.00和21.75小时，多项目跨室率25.04\%。

需要强调，59.53\%是可达折中点，不是“满足全部门诊、体检需求”的同义词。若医院政策规定更高的最低背景率，应在前沿上选择相应点，并接受住院率下降；若要求背景事件100\%在同日完成，则在当前标准班次、项目能力和医生容量下不可行，应增加资源、延长时段或把部分背景预约移日。

\begin{table}[H]
\centering
\caption{纯住院与推荐联合方案的运行形态}
\label{tab:scenario-operation}
\begin{tabular}{lrr}
\toprule
运行指标 & P2纯住院 & P3推荐联合方案 \\
\midrule
排入检查分钟 & 1,344,765 & 1,897,045 \\
人员在岗容量利用率 & 54.53\% & 76.93\% \\
住院等待P95/小时 & 23.67 & 25.58 \\
多项目需求跨室率 & 58.16\% & 30.63\% \\
背景全期下限/完成数 & -- & 43,221/54,137 \\
\bottomrule
\end{tabular}

\end{table}

表\ref{tab:scenario-operation}揭示了联合情景下降的运行原因。加入门诊、体检后，排入分钟增加53.72\%，人员在岗容量利用率由54.51\%升至83.79\%，住院等待P95由23.67小时增至38.58小时，已接近48小时边界。另一方面，住院多项目跨室率下降到25.04\%，并非联合方案改善了设备兼容，而是背景目标先占用部分灵活容量后，只有更容易形成完整路径的住院事件被排入。故跨室率下降不能单独当作服务改善指标。

$\alpha=0.30$的逐日目标覆盖364个有背景需求的日期，其中363日达到当天目标，仅1日短缺7个目标事件；全年最终背景完成57\,177例，远高于逐日预留阶段的28\,970例。这说明逐日配额只是构造算法的容量引导，最终可行性仍由全年$R_B\ge\alpha$判定。对$\alpha=0.65$和0.70，算法分别达到64.80\%和65.19\%，未达到所设下限；这两点用于标识可达边界，不能作为满足相应服务政策的方案。

\subsubsection{现实校准与瓶颈分解}

只报告最终60.38\%无法判断下降来自数据口径、设备能力还是算法。本文构造C0--C7校准链：

\begin{enumerate}
\item C0：按真实开单至报告提交计算历史留出期48小时率；
\item C1：只使用题给时长上界，不限制设备和标准班次；
\item C2：加入08:00--12:00、13:00--17:00班次和22并发上限；
\item C3：加入训练期星期—时隙医生容量；
\item C4：加入表1项目—设备兼容关系；
\item C5：加入项目级空腹、憋尿和床旁语义；
\item C6：加入门诊、体检背景竞争；
\item C7：使用正式确定性构造策略。
\end{enumerate}

\begin{table}[H]
\centering
\caption{历史现实校准与约束增量}
\label{tab:calibration}
\resizebox{\textwidth}{!}{\begin{tabular}{lrrr}
\toprule
校准层 & 完成数/分母 & 完成率 & 较上一层/百分点 \\
\midrule
历史开单至报告 & 59,266/70,864 & 83.63\% & -- \\
仅题给时长 & 70,771/70,771 & 100.00\% & +16.37 \\
加入标准班次 & 70,771/70,771 & 100.00\% & +0.00 \\
加入医生并发 & 70,771/70,771 & 100.00\% & +0.00 \\
加入项目设备能力 & 60,526/70,771 & 85.52\% & -14.48 \\
加入项目级准备 & 60,359/70,771 & 85.29\% & -0.24 \\
加入门诊/体检竞争 & 42,730/70,771 & 60.38\% & -24.91 \\
正式构造策略 & 42,730/70,771 & 60.38\% & +0.00 \\
\bottomrule
\end{tabular}
}
\end{table}

\begin{figure}[H]
\centering
\includegraphics[width=\textwidth]{fig06_calibration.pdf}
\caption{从历史口径到正式联合排程的逐层现实校准}
\label{fig:calibration}
\end{figure}

历史开单至报告率为83.63\%，但其报告时间包含医院真实加班、排班和流程，且事件口径与排程代理不同，不能直接视为模型基准。C1--C3均为100\%，表明若设备同质，医生总并发在标准班次内足以容纳纯住院题给工时；加入项目—设备兼容后降至85.52\%，说明专项设备集合是纯住院的首要限制。项目级准备使其再降0.24个百分点。加入门诊、体检竞争后下降24.91个百分点，说明联合场景的主要矛盾是共享容量分配，而非FCFS与松弛度的细微差别。

\subsubsection{推荐排程的运行形态}

推荐联合方案共导出176\,739条已排任务，涉及100\,303个完整或部分可执行事件；评价时只有完整事件进入完成分子。图\ref{fig:daily}显示日级率受周末、节假日和项目组合影响，14日滚动后仍存在明显波动，说明不能以总体率替代每日运营监控。

\begin{figure}[H]
\centering
\includegraphics[width=\textwidth]{fig08_daily_performance.pdf}
\caption{推荐联合策略的日级稳定性}
\label{fig:daily}
\end{figure}

图\ref{fig:gantt}给出总排入分钟最高的代表日设备甘特图。每个条块对应一个患者项目，资源结构使用题给22台设备、上午/下午工作块和专项兼容关系。午休时段为空白，床旁7号机和专项诊室的负荷由实际可行任务产生。

\begin{figure}[H]
\centering
\includegraphics[width=\textwidth]{fig09_representative_gantt.pdf}
\caption{高负荷代表日的设备排程}
\label{fig:gantt}
\end{figure}
